\section{The Same Measurements on a Second Workload and a Second Accelerator}
\label{sec:xtrace}
\label{sec:a100}

We repeated the same accounting, sweep and controller on a second replay---117
real Claude agent sessions, 390 rounds, same model, engine and
101{,}432-token pool---and the same sweep on an A100 (80\,GB) replaying the main
trace (nominal pool 983{,}728 tokens, 9.70$\times$ the 3090's). It ran
entirely without a host tier---device-only points, envelope and controller---so it
tests the negative result under a configuration that \emph{cannot} separate the
two.

\paragraph{The direction transfers; the level and the place do not.} Higher
planned reuse (86.3\% of prompt tokens against 75.5\%), less of it realized
(74.8\% to 57.5\% as the count rises, against 92.5\% to 36.4\%): the direction of
\S\ref{sec:model} survives, the sharpness does not, and we do not lean on slopes
spanning ranges of different length and level. Nor does the \emph{level} transfer:
realization is \emph{less} at 15.8\% median device occupancy than the main trace's
at a comparable count.

\paragraph{The tier mechanism moves with the tier.} The A100 criterion was fixed in
advance on one decision axis: the $\overline{R}_{\text{busy}}$ at which realized
reuse crosses 50\%. On the 3090 that crossing is 1.7344; the A100 never crosses---at
3.4077, twice the 3090's, it is still at 92.5\%. The pre-registered rule reads a
crossing, not a slope, and returns \emph{supported}. The same measurement in
$\epsilon$ is flat, not steep---$-0.74$ on the first segment against the 3090's
$-0.21$, then $-0.27$, $+0.03$, $+0.34$, $+0.72$, staying under one and reaching $+1.49$
on the last, a segment spanning 3.22 to 3.41 that carries the same
small-denominator caveat we flag on the 3090: read it as the crossing, not as a
rate. So added concurrency stops paying at roughly twice the running count on the
larger-pool card---not the pool's factor of $9.7$, and we do not claim it is; a
supporting measurement, not a second criterion.

The best controller of \S\ref{sec:results} there completes the replay---999 of 999
steps, no failures---in 88.3 minutes at a 100.0\% deadline satisfaction rate. The
sweep has since
extended to caps eight and sixteen: cap eight finishes in 90.1 minutes, cap sixteen
in 88.3---the same wall clock as the controller---with the front interpolated there
0.1\,pp below the controller's SLO, inside the $\pm1.0$\,pp band we judge in. The
verdict is therefore a tie, obtained by \emph{reaching} the controller's operating
point rather than stopping short of it. \emph{It is a weaker claim on this
evidence}: \emph{one}
controller run, our pre-registered criterion for this tier has one axis and does not
cover it, and the controller was rarely binding---its pool peaked at $52\%$
utilisation and it admitted after a median wait of 1.6\,s---so a tie is close to what
``no admission decision at all'' would give. What transfers is the mechanism; we do
not read the tie as the negative result reproducing at full strength.

\begin{table}[t]
  \centering
  \small
  \setlength{\tabcolsep}{4pt}%
  \begin{tabular}{llrrrrr}
    \toprule
    Machine & Pool & $\overline{R}_{\text{busy}}$ & Occ. & Reuse & $\tau_2$ & SLO \\
    \midrule
    3090-24G & 101{,}432  & 0.9980 & 22.4\% & 92.5\% & 20.6 & 94.8 \\
    A100-80G & 971{,}456  & 0.9997 &  2.1\% & 92.5\% & 14.3 & 99.9 \\
    \midrule
    3090-24G & 101{,}432  & 1.6505 & 66.4\% & 60.7\% & 25.8 & 55.7 \\
    A100-80G & 983{,}728  & 1.9702 &  6.5\% & 92.5\% & 14.5 & 100.0 \\
    \midrule
    3090-24G & 101{,}432  & 1.8409 & 74.3\% & 36.4\% & 31.2 & 27.1 \\
    A100-80G & 983{,}728  & 2.3356 &  8.0\% & 92.5\% & 15.5 & 100.0 \\
    \bottomrule
  \end{tabular}
  \caption{The same static caps on two accelerators, rows paired at the
  same nominal cap. The A100's cap-one row came up with a pool of 971{,}456
  rather than 983{,}728 (\S\ref{sec:a100:limits}), which does not affect its
  column. Occupancy is a share of the pool, $\tau_2$ is in milliseconds, SLO is
  the share of first tokens within $10$\,s, and reuse is \texttt{reuse\_tok} over
  the trace's own ceiling.}
  \label{tab:a100}
\end{table}

\subsection{What this section does and does not buy}
\label{sec:a100:limits}

This buys a result that is not one trace family's: the same controller still does
not beat the same kind of static cap on a second workload, and the tier mechanism
behaves as the premise predicts on a second accelerator. It costs us the ability
to \emph{locate} the criterion's turn---the $\pm10\%$ perturbation puts every
segment's bounds across zero on both traces---and the controller sits below both
turns; the second workload is closer to the regime where adaptive control
\emph{should} pay, and it still did not win there. Three things bound the hardware
half: three of the seven A100 runs did not hold everything but the card fixed---cap
one's pool 971{,}456 rather than 983{,}728, cap eight 943{,}541, cap sixteen
983{,}697---none pool-bound, so the shortfall does not reach the realization
column; our no-control baseline is not clean---a foreign process in 82 of the
run's 341 in-window samples, all of them in its last 23.8 minutes---so we report
it as indicative, and co-tenancy inflates a baseline, so the direction is
conservative; that run finished in 88.6 minutes at 99.7\% SLO against the
controller's 88.3; and the clean re-run we attempted afterwards did not produce
one: pool 817{,}527 rather than 983{,}728, a median occupancy of 95.3\%, the
watcher records a foreign process in 566 of 566 in-window samples,
148.1 minutes at a 32.9\% deadline satisfaction rate, and a per-request decode
interval of
222.7\,ms---a reading of a contended host, not a baseline, which we do not use
for anything.
